\documentclass[main.tex]{subfiles}
\begin{document}
\ifSubfilesClassLoaded{\begin{center}\small\itshape Part of the GnssProcessor Mathematical Reference \textbullet\ \today\end{center}}{}
\section{Terrestrial Reference Frames, Epochs, and Datum Transformation}\label{sec:trf}
% ============================================================
%  Frame vs epoch; Helmert similarity and its rate field;
%  ITRF/NAD83/NATRF2022 parameter sets; the NGS SIF occupation
%  model; a fully numerical CORV worked example.
% ============================================================

% --- Doc-local symbols ---------------------------------------------------
% NOTE: gnssreport.sty is the single source of truth for shared notation.
% Everything below is scoped to this component; nothing here is global,
% and no shared macro (\rxpos, \Tr, \dt, \los, ...) is redefined.
\newcommand{\helmT}{\bm{T}}                      % Helmert translation (m)
\newcommand{\helmE}{\bm{E}}                      % linearised rotation matrix
\newcommand{\stavel}{\bm{v}}                     % station velocity (m/yr)
\newcommand{\eulerv}{\bm{\omega}}                % Euler (plate rotation) vector
\newcommand{\epo}{t}                             % coordinate epoch
\newcommand{\refep}{t_0}                         % parameter reference epoch

A terrestrial coordinate is three numbers, a frame, and an epoch.  The
engine works in ITRF at the epoch of observation, so the frame and the
epoch are usually implicit.  They become explicit when a result is
compared with a national datum, a published datasheet, or a passive
mark.  Changing the frame and changing the epoch are different
operations, and mixing them up produces decimetre-level offsets.  The
rule is:

\begin{empheq}[box=\fbox]{equation}
\textbf{Frames take a Helmert. Epochs take a velocity.}
\label{eq:trf-rule}
\end{empheq}

\Cref{sec:trf-rule,sec:trf-veltrans,sec:trf-plate} develop the rule and
its one exception.  \Cref{sec:trf-conv,sec:trf-params} give the sign
conventions and the published parameter sets.  \Cref{sec:trf-natrf}
covers NATRF2022.  \Cref{sec:trf-sif,sec:trf-arp} cover the NGS Site
Information File and the monument-to-ARP offset.  \Cref{sec:trf-corv}
is a numerical example at CORV, and \Cref{sec:trf-defmodel} compares the
NGS deformation model with station velocities.

% ============================================================
\subsection{Frame versus epoch}\label{sec:trf-rule}

A Helmert similarity (seven parameters, or fourteen with their linear
rates) maps frame $A$ to frame $B$ \emph{at the same epoch}.  It
describes how two realisations differ in origin, orientation, and
scale.  It says nothing about how the ground moves.

The fourteen-parameter form is
\begin{equation}
p(\epo) \;=\; p_0 \;+\; \dot{p}\,(\epo-\refep),
\qquad p \in \{\helmT,\,s,\,\bm{\varepsilon}\},
\label{eq:trf-14p}
\end{equation}
and the rates $\dot{p}$ are rates \emph{of the parameters}: the drift of
one frame relative to the other.  They are not station velocities.
Evaluating \eqref{eq:trf-14p} at epoch $\epo$ gives the transformation
from $A$ to $B$ at that epoch.  It does not move the coordinate through
time.

Moving a coordinate between epochs requires information about the site:
\begin{equation}
\bm{x}(\epo) \;=\; \bm{x}(\refep) \;+\; \stavel\,(\epo-\refep)
              \;+\; \bm{d}_{\mathrm{per}}(\epo)
              \;+\; \bm{d}_{\mathrm{PSD}}(\epo),
\label{eq:trf-propagate}
\end{equation}
with $\stavel$ a station velocity, $\bm{d}_{\mathrm{per}}$ the periodic
(annual $+$ semiannual) terms, and $\bm{d}_{\mathrm{PSD}}$ the
post-seismic deformation.  \texttt{resolveStationCoordinate} and
\texttt{tools/sif\_propagate} implement this form.  The three site terms
are observed or modelled.  A transformation parameter set does not
contain them, so a transformation that appears to handle an epoch
change is using the rigid-plate approximation of \Cref{sec:trf-plate}.


\begin{figure}[htbp]
\centering
\begin{tikzpicture}[node distance=14mm and 30mm,
  fr/.style={proc, minimum width=31mm, minimum height=9mm},
  lbl/.style={font=\footnotesize, fill=white, inner sep=1.5pt}]
  % rows = frames, columns = epochs
  \node[fr, very thick]           (N10) {NAD83(2011)\\{\footnotesize 2010.00}};
  \node[fr, right=of N10]         (N20) {NAD83(2011)\\{\footnotesize 2020.0}};
  \node[fr, below=of N10]         (I10) {ITRF2020\\{\footnotesize 2010.0}};
  \node[fr, below=of N20, very thick] (I20) {ITRF2020\\{\footnotesize 2020.0}};
  \node[fr, right=of I20]         (It)  {ITRF2020\\{\footnotesize $t$}};
  \node[fr, below=of I20]         (T20) {NATRF2022\\{\footnotesize 2020.00}};
  \node[fr, below=of It]          (Tt)  {NATRF2022\\{\footnotesize $t$}};
  % epoch changes: velocity (horizontal)
  \draw[flow, <->, very thick] (I10) -- node[lbl] {$\stavel$} (I20);
  \draw[flow, <->] (I20) -- node[lbl] {$\stavel$} (It);
  \draw[flow, <->] (N10) -- node[lbl] {$\stavel_{\mathrm{NAD83}}$} (N20);
  \draw[flow, <->] (T20) -- node[lbl] {$\stavel-\eulerv\times\bm{x}$} (Tt);
  % frame changes: Helmert / Euler pole (vertical)
  \draw[flow, <->, very thick] (I10) -- node[lbl] {Helmert(2010.0)} (N10);
  \draw[flow, <->] (I20) -- node[lbl] {Helmert(2020.0)} (N20);
  \draw[flow, <->] (It)  -- node[lbl] {Euler pole $(t-2020)$} (Tt);
  \draw[double, double distance=1.5pt, semithick] (I20) -- node[lbl] {$\equiv$} (T20);
  % axes
  \node[note, above=3mm of N10, xshift=31mm] {epoch change $\longleftrightarrow$ station velocity};
  \node[note, rotate=90, anchor=south] at ([xshift=-5mm]$(N10.west)!0.5!(I10.west)$) {frame change $\updownarrow$ Helmert};
\end{tikzpicture}
\caption{Frames (rows) and epochs (columns).  Vertical moves change the
frame at a fixed epoch; horizontal moves change the epoch within a frame
and need a station velocity.  Heavy lines mark the worked example of
\Cref{sec:trf-corv}, from ITRF2020 at 2020.0 to NAD83(2011) at 2010.00.
Because NAD83 rotates with the plate, $\stavel_{\mathrm{NAD83}}$ is
small on the stable interior and larger at the plate margins
(\Cref{sec:trf-plate}).}
\label{fig:trf-grid}
\end{figure}

% ============================================================
\subsection{Transforming a velocity}\label{sec:trf-veltrans}

Differentiating the linearised transformation with respect to time gives
\begin{equation}
\stavel_B \;=\; \stavel_A \;+\; \dot{\helmT}
              \;+\; \dot{s}\,\bm{x} \;+\; \dot{\helmE}\,\bm{x}.
\label{eq:trf-veltrans}
\end{equation}
Three properties of \eqref{eq:trf-veltrans} matter in practice.

\paragraph{The rate field is rigid.}  The added terms
$\dot{\helmT} + \dot{s}\bm{x} + \dot{\helmE}\bm{x}$ depend on the
station's \emph{position} only, not on its motion.  This was checked by
transforming fictitious stations at one position with different
velocities: the added velocity was identical to the micron.  The rates
of a fourteen-parameter set are therefore a rigid velocity field, just
as the seven parameters are a rigid displacement field.

\paragraph{Propagation and transformation commute, with one condition.}
Propagating in frame $A$ and then transforming gives the same result as
transforming and then propagating in frame $B$, to $0.000$~mm, provided
each propagation uses the velocity expressed in its own frame.  That
velocity comes from \eqref{eq:trf-veltrans}.

\paragraph{Using the wrong frame's velocity is a large error.}
Propagating a NAD83(2011) coordinate with an ITRF velocity costs about
$20$~mm/yr at CORV, or $203$~mm per decade.  No single-epoch check
reveals it.  The error exceeds every geophysical correction of
\dcref{sec:sitedisp}{the Site Displacement document} combined.

% ============================================================
\subsection{The rigid-plate approximation}\label{sec:trf-plate}

The parameter rates can stand in for a station velocity in one case.
A plate-fixed frame is built so that stations on the rigid plate do not
move in it.  Setting $\stavel_B = \bm{0}$ in \eqref{eq:trf-veltrans}
gives
\begin{equation}
\stavel_A \;=\; -\bigl(\dot{\helmT} + \dot{s}\,\bm{x}
                        + \dot{\helmE}\,\bm{x}\bigr).
\label{eq:trf-plate}
\end{equation}
The right-hand side is then a \emph{modelled} rigid-plate velocity, and
\eqref{eq:trf-14p} is being used as a one-term deformation model.

The approximation is as good as the assumption.  On stable North
America the horizontal velocity relative to the plate is
$1$--$2$~mm/yr, so \eqref{eq:trf-plate} costs $1$--$3$~cm per decade.
At a plate margin it fails.  CORV sits on the deforming Cascadia
margin, and the fourteen-parameter-only reconstruction of
\Cref{sec:trf-corv} misses by $104$~mm per decade.


% ============================================================
\subsection{Rotation-sign conventions}\label{sec:trf-conv}

Two rotation conventions are in use, and the sources cited here use
both.  The NGS / EPSG:1056 \emph{coordinate frame} form, linearised for
small rotations, is
\begin{equation}
\bm{x}' \;=\; \bm{x} \;+\; \helmT \;+\; s\,\bm{x} \;+\; \helmE\,\bm{x},
\qquad
\helmE \;=\;
\begin{bmatrix}
 0        & \varepsilon_z & -\varepsilon_y \\
-\varepsilon_z & 0        &  \varepsilon_x \\
 \varepsilon_y & -\varepsilon_x & 0
\end{bmatrix}.
\label{eq:trf-cf}
\end{equation}
The IERS / IGN \emph{position vector} form is identical except that
$\helmE$ carries the opposite sign throughout; equivalently,
$\bm{\varepsilon} \to -\bm{\varepsilon}$.

Using one convention's parameters in the other convention's matrix
applies the rotation backwards.  For a rotation angle $\theta$ and a
position of magnitude $L$ the error is $2\lvert\theta\rvert L$, twice
the size of omitting the rotation.  The same holds for the rate field
of \Cref{sec:trf-veltrans}.  A rate field derived here with the wrong
convention came out exactly negated, a $\sim\!35$~mm/yr error at CORV
where the true field is $\sim\!20$~mm/yr.  Check the sign of a rate
field against a station of known velocity; a plausible magnitude is not
evidence of a correct sign.

% ============================================================
\subsection{Published parameter sets and their composition}\label{sec:trf-params}

\Cref{tab:trf-params} lists the three sets the tooling uses.  The first
is in position-vector convention and the other two are in
coordinate-frame convention.  The IGN rotations are zero, so the
composition below can add the sets without converting between the
conventions of \Cref{sec:trf-conv}.

\begin{table}[htbp]
\centering
\small
\caption{Published Helmert parameter sets.  ITRF2020$\to$ITRF2014 is
IGN's, in position-vector convention; the two NAD83 sets are NGS/EPSG,
in coordinate-frame convention.  Rotations $\bm{\varepsilon}$ in
milliarcseconds.}
\label{tab:trf-params}
\begin{tabular}{llll}
\toprule
& \textbf{ITRF2020$\to$ITRF2014} & \textbf{ITRF2014$\to$NAD83(2011)} & \textbf{ITRF2020$\to$NAD83(2011)} \\
& IGN, position-vector & NGS MYCS2, EPSG:8970 & EPSG:10334 (pre-composed) \\
& $\refep=2015.0$ & $\refep=2010.0$, coord-frame & $\refep=2010.0$, coord-frame \\
\midrule
$T_x$ & $-1.4$ mm            & $1.00530$ m      & $1.00390$ m \\
$T_y$ & $-0.9$ mm            & $-1.90921$ m     & $-1.90961$ m \\
$T_z$ & $+1.4$ mm            & $-0.54157$ m     & $-0.54117$ m \\
$s$   & $-0.42$ ppb          & $0.36891$ ppb    & $-0.05109$ ppb \\
$\varepsilon_x$ & $0$        & $26.78138$ mas   & $26.78138$ mas \\
$\varepsilon_y$ & $0$        & $-0.42027$ mas   & $-0.42027$ mas \\
$\varepsilon_z$ & $0$        & $10.93206$ mas   & $10.93206$ mas \\
\midrule
$\dot{T}_x$ & $0.0$ mm/yr    & $0.00079$ m/yr   & $0.00079$ m/yr \\
$\dot{T}_y$ & $-0.1$ mm/yr   & $-0.00060$ m/yr  & $-0.00070$ m/yr \\
$\dot{T}_z$ & $+0.2$ mm/yr   & $-0.00144$ m/yr  & $-0.00124$ m/yr \\
$\dot{s}$   & $0$            & $-0.07201$ ppb/yr& $-0.07201$ ppb/yr \\
$\dot{\varepsilon}_x$ & $0$  & $0.06667$ mas/yr & $0.06667$ mas/yr \\
$\dot{\varepsilon}_y$ & $0$  & $-0.75744$ mas/yr& $-0.75744$ mas/yr \\
$\dot{\varepsilon}_z$ & $0$  & $-0.05133$ mas/yr& $-0.05133$ mas/yr \\
\bottomrule
\end{tabular}
\end{table}

PROJ applies the pre-composed EPSG:10334 set for
\texttt{EPSG:9988}$\to$\texttt{EPSG:6317}, with
\texttt{t\_epoch}$=2010$ and the \texttt{coordinate\_frame} convention.
It can be rebuilt from the other two sets in two steps:

\begin{enumerate}[leftmargin=*]
\item \textbf{Run the IGN set to the NGS epoch.}  Evaluate
      \eqref{eq:trf-14p} at $2010.0$ from its own $\refep=2015.0$:
      $\helmT(2010) = \helmT(2015) + \dot{\helmT}\cdot(-5)
       = (-1.4,\,-0.4,\,+0.4)$~mm.
\item \textbf{Add term by term} to the NGS set at its native $2010.0$.
\end{enumerate}

The result matches PROJ's set: the largest difference over the fourteen
parameters is $4\times10^{-15}$ in the respective units.  Adding the
sets is valid because the cross terms of the exact composition are
below a micron at Earth radius.  The rotation-times-translation term
$\bm{\Omega}\helmT$ is $\sim\!0.3$~\si{\micro\metre} and the
rotation-squared term $\Omega_1\Omega_2 R$ is
$\sim\!0.1$~\si{\micro\metre}.  The two-step chain and the pre-composed
set agree to $0.0000$~mm at $2010.0$, $2015.0$, and $2020.0$.

Two conditions apply.  First, both sets must be in the same rotation
convention before they are added.  That step is unnecessary here only
because the IGN rotations are zero.  Second, step~1 is required: adding
the IGN set at its own $2015.0$ reference epoch introduces a $0.5$~mm
error in $Y$ and $1.0$~mm in $Z$.

The EPSG codes are ITRF2020 \texttt{9988} /
\texttt{9990} / \texttt{9989} (geocentric XYZ, 2D, 3D), ITRF2014
\texttt{7789} / \texttt{9000} / \texttt{7912}, and NAD83(2011)
\texttt{6317} / \texttt{6318} / \texttt{6319}.

% ============================================================
\subsection{NATRF2022}\label{sec:trf-natrf}

NATRF2022, the North American frame of NSRS2022, is ITRF2020 rotated
about the North America Euler pole.  It has no translation and no
scale.  The two frames are identical at epoch $2020.00$ and separate as
the plate moves.  With $\eulerv$ the Euler vector, positive
counter-clockwise under the right-hand rule,
\begin{align}
\bm{x}_{\mathrm{NATRF}}(\epo)
  &= \bm{x}_{\mathrm{ITRF2020}}(\epo)
     \;-\; (\eulerv \times \bm{x})\,(\epo - 2020.0),
     \label{eq:trf-natrf-pos}\\
\stavel_{\mathrm{NATRF}} &= \stavel_{\mathrm{ITRF2020}} - \eulerv \times \bm{x}.
     \label{eq:trf-natrf-vel}
\end{align}

The NGS \textsc{beta} EPP2022 Euler poles are, in mas/yr,
$\eulerv_{\mathrm{NA}} = (0.046,\,-0.704,\,-0.047)$, with the companion
plates $\eulerv_{\mathrm{PA}} = (-0.409,\,1.063,\,-2.188)$,
$\eulerv_{\mathrm{CA}} = (-0.056,\,-0.957,\,0.589)$, and
$\eulerv_{\mathrm{MA}} = (-8.089,\,5.937,\,2.159)$ for PATRF2022,
CATRF2022, and MATRF2022 respectively.  These are \textsc{beta} values
from \url{https://beta.ngs.noaa.gov/NATRF2022/}.  Fetch the current
values before production use.

The sign of \eqref{eq:trf-natrf-vel} was checked on stable-interior
stations, whose plate-fixed velocity must be near zero.  The NATRF2022
velocities, ENU in mm/yr, are AMC4 $(-0.49,\,1.01,\,-1.17)$,
GODN $(-0.12,\,0.47,\,-1.49)$, TXAU $(-0.34,\,1.73,\,-0.52)$,
WES2 $(0.07,\,-0.03,\,-0.12)$, and MIA3 $(0.28,\,0.90,\,0.56)$.  A sign
error would double the plate rate, giving $20$--$30$~mm/yr here.

PROJ 9.5.1 has no NATRF2022 CRS, so the caller must apply
\eqref{eq:trf-natrf-pos} directly.

% ============================================================
\subsection{SIF occupations and the datasheet convention}\label{sec:trf-sif}

The NGS Site Information File (SIF) gives ITRF2020 coordinates at epoch
$2020.0$ from the MYCS3 solution.  A station can have several
\textbf{occupation} windows, each with its own reference coordinate and
velocity.  An occupation boundary is an equipment change or a known
offset, so the coordinate series is discontinuous there.  Two
consequences follow.

\paragraph{The datasheet convention.}  NGS produces a datasheet's
NAD83(2011) epoch-$2010.00$ coordinate as follows.  Take the
\emph{latest} occupation.  Propagate its reference coordinate back ten
years with its linear velocity and \textbf{no} periodic term.  Apply
the fourteen-parameter transformation of \Cref{sec:trf-params} at
$2010.0$.  This recipe reproduces the datasheet to $0.4$~mm at CORV
(PID AH2488) and $0.6$~mm at CHZZ (PID AJ6961).  In
\texttt{tools/sif\_propagate} it is the \texttt{-{}-datasheet} option.

\paragraph{The engine convention.}  \texttt{resolveStationCoordinate},
and \texttt{tools/sif\_propagate} by default, select the occupation that
\emph{covers the requested epoch}.  A coordinate for an observation
session must describe the equipment installed at that time.  At CORV a
request for $2010.0$ uses the 1998--2019 occupation and differs from the
datasheet by $32.8$~mm.  Of that, $\delta U = -32.4$~mm is the jump at
the 2019-07-24 antenna change (ASH700936E\_C $\to$ TRM59800.99 SCIS).
Both coordinates are correct for their purpose.  One rule follows:

\begin{remark}
The difference of two resolved coordinates across an occupation break
is \emph{not} a velocity.  It includes the equipment offset, which is
often the larger part.
\end{remark}

\paragraph{Periodic terms are a constant offset at whole-year epochs.}
The SIF annual and semiannual terms are referenced to the station's
reference epoch.  A whole number of years from that epoch,
$\cos = 1$ and $\sin = 0$, so the terms add a fixed vector: $0.65$~mm at
CORV, $1.73$~mm at NEWP, $3.18$~mm at CHZZ.  The offset cancels in a
difference between two such epochs but not in an absolute comparison.
Published coordinates exclude it, so a propagated coordinate that
includes it differs from the datasheet by these amounts.

% ============================================================
\subsection{Monument versus ARP}\label{sec:trf-arp}

A datasheet comparison must also use the same reference point.  SIF and
datasheet coordinates refer to the \textbf{monument} (ground reference
point, GRP).  The engine's position state and the NGS CORS coordinate
files refer to the \textbf{ARP}.  The offset between them is the RINEX
\texttt{ANTENNA: DELTA H/E/N} eccentricity, which the SIF publishes as
\texttt{ARP\_MINUS\_GRP\_ENU}.  The reduction is developed in
\dcref{sec:frame-arp}{the marker/ARP section of the Coordinate Frames
document}.

At CORV, $\delta H = 1.521$~m, and both comparisons close:
\begin{itemize}[leftmargin=*]
\item SIF monument $+\,1.521$~m Up versus the published ITRF2020 ARP
      $(-2498424.007,\,-3802820.797,\,4454737.751)$: $0.3$~mm.
\item Datasheet $+\,1.521$~m Up versus the published NAD83(2011) ARP
      $(-2498423.096,\,-3802822.020,\,4454737.741)$: $0.5$~mm.
\end{itemize}
Omitting this step at CORV gives an obvious $1.5$~m height error.  The
harder case is a site with an eccentricity of a few centimetres, where
the error is small enough to go unnoticed.

% ============================================================
\subsection{Worked example: CORV}\label{sec:trf-corv}

This example works every step above with numbers, at the monument,
using the latest occupation (2019-08-04 to open).
\texttt{tools/sif\_propagate/corv\_worked\_example.py} reproduces it.
% The reproducer reads the IFDM2022 grids from data/static/ifdm2022/
% (git-ignored; downloaded from the beta NCAT Downloads tab).


\begin{figure}[htbp]
\centering
\begin{tikzpicture}[node distance=15mm and 9mm,
  st/.style={io, text width=70mm, minimum height=11mm},
  op/.style={font=\footnotesize, align=left, anchor=west},
  chk/.style={proc, text width=38mm, minimum height=9mm}]
  \node[st]              (S1) {\textbf{ITRF2020}, 2020.0 \hfill{\footnotesize SIF}\\
        {\footnotesize $-2498423.4121\;\;{-}3802819.8917\;\;4454736.6830$}};
  \node[st, below=of S1] (S2) {\textbf{ITRF2020}, 2010.0\\
        {\footnotesize $-2498423.3133\;\;{-}3802819.9155\;\;4454736.7162$}};
  \node[st, below=of S2] (S3) {\textbf{ITRF2014}, 2010.0\\
        {\footnotesize $-2498423.3136\;\;{-}3802819.9143\;\;4454736.7147$}};
  \node[st, below=of S3] (S4) {\textbf{NAD83(2011)}, 2010.00\\
        {\footnotesize $-2498422.5017\;\;{-}3802821.1141\;\;4454736.6736$}};
  \draw[flow] (S1) -- node[op, xshift=2mm] {velocity $\times\,(-10$ yr)\\$(-9.88,\,+2.38,\,-3.31)$ mm/yr} (S2);
  \draw[flow] (S2) -- node[op, xshift=2mm] {Helmert at 2010.0\\$\helmT=(-1.4,-0.4,+0.4)$ mm,\; $s=-0.42$ ppb} (S3);
  \draw[flow] (S3) -- node[op, xshift=2mm] {Helmert at $\refep=2010.0$\\fourteen-parameter set, \Cref{sec:trf-params}} (S4);
  % checks
  \node[chk, right=of S4]  (P) {PROJ \texttt{EPSG:10334}\\{\footnotesize $0.0$ mm}};
  \node[chk, below=4mm of P] (D) {NGS datasheet AH2488\\{\footnotesize $0.4$ mm}};
  \draw[feedback] (S4.east) -- (P.west);
  \draw[feedback] (S4.east) -- ++(4mm,0) |- (D.west);
  % monument -> ARP
  \node[chk, right=of S1] (A1) {$+\,1.521$ m Up $=$ ARP\\{\footnotesize NGS CORS file: $0.3$ mm}};
  \node[chk, below=4mm of D] (A4) {$+\,1.521$ m Up $=$ ARP\\{\footnotesize NGS CORS file: $0.5$ mm}};
  \draw[feedback] (S1) -- (A1);
  \draw[feedback] (S4.east) -- ++(4mm,0) |- (A4.west);
\end{tikzpicture}
\caption{The CORV worked example, modern $\to$ legacy, at the monument
(metres).  Solid arrows are the computation; dashed arrows are
independent checks, with the 3-D closure.  The antenna reference point
lies $1.521$~m above the monument (RINEX \texttt{DELTA H}).}
\label{fig:trf-corv}
\end{figure}

\paragraph{Starting point (SIF, ITRF2020 @ 2020.0).}
$\bm{x} = (-2498423.4121,\,-3802819.8917,\,4454736.6830)$~m,
$h = 105.5881$~m; velocity in ECEF
$(-9.8836,\,2.3812,\,-3.3134)$~mm/yr, i.e.\ ENU
$(-9.57,\,-4.77,\,+0.12)$~mm/yr.

\paragraph{Step 2a: propagate to 2010.0 in ITRF2020.}
$\bm{x} - 10\stavel = (-2498423.3133,\,-3802819.9155,\,4454736.7162)$~m;
the station moved $106.9$~mm.

\paragraph{Step 2b: ITRF2020 $\to$ ITRF2014 at 2010.0.}
With $\helmT = (-1.4,\,-0.4,\,+0.4)$~mm (the IGN set run to $2010.0$)
and $s = -0.42$~ppb:
$(-2498423.3136,\,-3802819.9143,\,4454736.7147)$~m.

\paragraph{Step 2c: ITRF2014 $\to$ NAD83(2011) at 2010.0.}
Translation $(1.00530,\,-1.90921,\,-0.54157)$~m; rotation term
$(-0.19247,\,0.71082,\,0.49885)$~m; scale term
$(-0.00092,\,-0.00140,\,0.00164)$~m.  Result
$(-2498422.5017,\,-3802821.1141,\,4454736.6736)$~m.  PROJ gives the
same result ($0.0$~mm).

\paragraph{Closure against the datasheet.}  NGS datasheet AH2488 gives
$\varphi = 44^\circ\,35'\,07.91300''$~N,
$\lambda = 123^\circ\,18'\,16.51681''$~W, $h = 105.953$~m, i.e.\
$(-2498422.5015,\,-3802821.1143,\,4454736.6734)$~m.  Hand computation
minus datasheet is $\delta E = -0.3$, $\delta N = +0.2$,
$\delta U = +0.1$~mm, which is $0.4$~mm in three dimensions.

\paragraph{The reverse direction.}  Going from the datasheet to ITRF2020
at $2020.00$, which is also NATRF2022 at that epoch
(\Cref{sec:trf-natrf}), requires ten years of station velocity.  The
miss against the SIF coordinate depends on the source of that velocity:

\begin{table}[htbp]
\centering\small
\caption{CORV reverse reconstruction: miss against the SIF coordinate by
velocity source.  The last row is the rigid-plate approximation of
\Cref{sec:trf-plate}, which does not hold at CORV.}
\label{tab:trf-corv-reverse}
\begin{tabular}{lll}
\toprule
\textbf{Velocity source} & \textbf{3D miss} & \textbf{Dominant component} \\
\midrule
SIF velocity                       & $0.4$~mm   & $\delta E = 0.3$ \\
IFDM2022 grid $(-8.68,-4.78,+0.04)$ mm/yr ENU & $9.3$~mm & $\delta E = 9.2$ \\
\textsc{beta} NCAT                 & $9.6$~mm   & $\delta E\,6.9$, $\delta N\,-5.4$, $\delta U\,4.0$ \\
14-parameter only ($\stavel_{\mathrm{NAD83}} := \bm{0}$) & $104.2$~mm & $\delta E\,-54.6$, $\delta N\,-88.2$, $\delta U\,10.4$ \\
\bottomrule
\end{tabular}
\end{table}

\paragraph{The same point in four frames.}  CORV velocity, ENU mm/yr:
ITRF2020 $(-9.57,\,-4.77,\,0.12)$; plate motion $\eulerv\times\bm{x}$
$(-13.20,\,-13.12,\,-0.04)$; NATRF2022 $(3.63,\,8.35,\,0.17)$;
NAD83(2011) $(5.49,\,8.80,\,-1.05)$.  The NATRF2022 velocity is what
remains after removing rigid North America.  At a deforming margin it
is not small: a NATRF2022 coordinate at CORV moves
$(21.8,\,50.1,\,1.0)$~mm by $2026.0$ and $(36.3,\,83.5,\,1.7)$~mm by
$2030.0$ from its $2020.00$ value.  A plate-fixed coordinate is not a
static one.


\begin{figure}[htbp]
\centering
\begin{tikzpicture}[node distance=5mm and 11mm,
  tool/.style={proc, text width=46mm, minimum height=10mm},
  res/.style={io, minimum width=17mm, minimum height=9mm},
  hd/.style={font=\small\bfseries, align=center}]
  \node[tool]                 (A) {PROJ $+$ \texttt{sif\_propagate}\\{\footnotesize station velocity (SIF)}};
  \node[tool, below=of A]     (B) {PROJ $+$ IFDM2022 grid\\{\footnotesize model velocity}};
  \node[tool, below=of B]     (C) {NCAT (beta)\\{\footnotesize model velocity (IFDM2022)}};
  \node[tool, below=of C]     (D) {PROJ only\\{\footnotesize no velocity (parameter rates)}};
  \node[io, left=of $(B.west)!0.5!(C.west)$, minimum height=14mm] (src) {NAD83(2011)\\2010.00};
  \foreach \n in {A,B,C,D} \draw[flow] (src.east) -- ++(4mm,0) |- (\n.west);
  \node[res, right=of A] (A1) {0.4};  \node[res, right=4mm of A1] (A2) {0.0};
  \node[res, right=of B] (B1) {9.3};  \node[res, right=4mm of B1] (B2) {11.8};
  \node[res, right=of C] (C1) {9.6};  \node[res, right=4mm of C1] (C2) {17.1};
  \node[res, right=of D, fill=black!12] (D1) {\textbf{104.2}}; \node[res, right=4mm of D1] (D2) {26.1};
  \foreach \n in {A,B,C,D} \draw[flow] (\n.east) -- (\n1.west);
  \node[hd, above=2mm of A1] {CORV\\{\footnotesize\mdseries margin}};
  \node[hd, above=2mm of A2] {AMC4\\{\footnotesize\mdseries interior}};
  \node[hd, above=2mm of A]  {ten years supplied by};
  \node[note, below=2mm of $(D1.south)!0.5!(D2.south)$, align=center] {miss at ITRF2020 $\equiv$ NATRF2022, 2020.00 [mm]};
\end{tikzpicture}
\caption{Legacy $\to$ modern, NAD83(2011) epoch 2010.00 to
ITRF2020 $\equiv$ NATRF2022 epoch 2020.00, by four routes.  PROJ is the
inverse of \texttt{EPSG:10334}, evaluated at 2010.0 (at 2020.0 in the
last row, where it stands alone).  The number is
the 3-D miss against the SIF coordinate at 2020.0.  CORV starts from NGS
datasheet AH2488; AMC4 (Colorado Springs, single occupation) starts from
its SIF-derived 2010.00 coordinate, so its first row is zero by
construction.}
\label{fig:trf-tools}
\end{figure}

% ============================================================
\subsection{Deformation models and validation}\label{sec:trf-defmodel}

A passive mark has no observed velocity, so a model must supply one.
NGS supplies IFDM2022 with \textsc{beta} NCAT.  It is a global
$0.1^\circ$ ENU velocity grid with per-node standard deviations, in
ITRF2020 (\texttt{epsg:9990}), plus 33 earthquake displacement grids.
It was compared with SIF velocities at 21 stations in CONUS, Alaska,
and British Columbia, using one occupation per station:

\begin{itemize}[leftmargin=*]
\item $\lVert \stavel_{\mathrm{IFDM}} - \stavel_{\mathrm{SIF}} \rVert$:
      median $1.34$~mm/yr, maximum $3.11$~mm/yr.
\item The grid's formal $\sigma$ is optimistic by roughly a factor of
      two (median ratio $2.1$, up to $5.9$ in the stable interior).
      The comparison is \emph{in-sample}, because the grid was built
      from these CORS.  A passive mark should expect a larger error.
\item In the stable interior the IFDM horizontal velocity equals the
      EPP2022 plate motion to $0.01$~mm/yr.  The model assigns zero
      intraplate velocity there.
\item Velocity spread across a single grid cell is $0.03$--$0.05$~mm/yr
      in the interior and up to $3.7$~mm/yr in the Aleutians.
      Interpolation error matters only where the field is steep.
\end{itemize}

NCAT's transformation from NAD83(2011) at $2010.00$ to NATRF2022 at
$2020.00$ was compared with the SIF coordinate at $2020.0$.  The mean
difference is $(+4.9,\,-6.1,\,-1.3)$~mm and the median 3D difference is
$18$~mm, within NCAT's stated uncertainty of about $2$~cm for most
stations.  The small mean confirms that NATRF2022 and ITRF2020 coincide
at $2020.00$ to a few millimetres, as \Cref{sec:trf-natrf} requires.
Five stations are outliers: FAIR $115$, SNI1 $82$, AB02 $78$, KYW1 $61$,
and HOLP $43$~mm.  Four of them lie inside IFDM earthquake displacement
grids.  The SIF represents those events as occupation \emph{breaks}
rather than PSD terms, so the two sources model them differently.  This
attribution rests on grid coverage only.  The displacement magnitudes
were not evaluated, and KYW1 is unexplained.

\paragraph{Limitations of \textsc{beta} NCAT.}  It offers only fixed
epochs: NAD83(2011) at $2010.00$ and NATRF2022 at $2020.00$.  Its
offline CLI performs coordinate \emph{conversions} only, not datum
transformations.  Its XYZ input is rounded to $1$~mm.  The post-seismic
terms of both tools are untested here, because no station in the
comparison has a PSD term.

\paragraph{Mixing velocity sources.}  Within NSRS2022 the deformation
model is the convention, so using IFDM (or HTDP) for a passive mark is
correct by definition.  The risk is mixing sources.  A network that
propagates its CORS with SIF velocities and its marks with grid
velocities introduces a $1$--$2$~mm/yr difference between the two
groups, which distorts the network over a decade.  The grid also gives
only a diagonal covariance, and an optimistic one by the ratio above.
SIF propagation gives the full
\begin{equation}
\bm{C}_{xx}(\epo) \;=\; \bm{C}_{xx}
   \;+\; \Delta\epo\,\bigl(\bm{C}_{xv} + \bm{C}_{xv}^{\Tr}\bigr)
   \;+\; \Delta\epo^{2}\,\bm{C}_{vv},
\label{eq:trf-covprop}
\end{equation}
including the position--velocity cross-covariance, which dominates
\eqref{eq:trf-covprop} at short $\Delta\epo$ and which the grid does not
provide.

\ifSubfilesClassLoaded{\printbibliography}{}
\end{document}
